\documentclass[../main.tex]{subfiles}

\begin{document}
\chapter{Introducción}
\label{C:introduccion}

\section{Motivación}
\label{S:motivacion}

En las ciencias de la computación, las tecnologías avanzan rápidamente, y los sistemas antiguos suelen quedar obsoletos; sin embargo, ocasionalmente surge la necesidad de recuperar el uso de máquinas históricas, motivados por el interés histórico o educativo.

Durante la segunda mitad de los años 80 y la primera de los 90, la industria del entretenimiento experimentó un notable crecimiento debido a la proliferación de fabricantes de hardware a nivel mundial. Ordenadores domésticos como el ZX Spectrum o el Amstrad CPC llegaron a los hogares europeos, mientras que videoconsolas de 8 bits, basadas en procesadores como el Z80 o el 68000 y procedentes de Asia, permitieron que muchos usuarios tuvieran sus primeros contactos con software desarrollado, como la Nintendo Entertainment System o la portátil Game Boy.

El desarrollo de software para estas plataformas requiere conocimientos específicos de sus arquitecturas y limitaciones. En aquella época, los recursos eran reducidos y la mayoría del software se programaba en ensamblador directamente para el procesador objetivo, sin herramientas modernas de desarrollo.

En la actualidad, aunque existen compiladores para estas máquinas, su uso requiere gestionar herramientas dispersas y un conocimiento avanzado del entorno, lo que limita la accesibilidad y aumenta la complejidad del desarrollo.

El desarrollo enfocado a estas plataformas, además tiene ventajas formativas de cara a la preparación de futuros desarrolladores ya que centra el desarrollo de software en la optimización de recursos disponibles y la programación cercana al hardware.

Otro aspecto relevante es la preservación del patrimonio cultural digital. Aunque actualmente puede encontrarse información sobre sistemas históricos en internet, esta se halla dispersa y depende en gran medida de iniciativas individuales y voluntarias. La falta de centralización y de mecanismos formales de conservación implica el riesgo de que, a medio o largo plazo, parte del conocimiento técnico asociado a estas arquitecturas, así como el software desarrollado para ellas, pueda desaparecer. En este sentido, el presente proyecto contribuye a la preservación de esta cultura tecnológica mediante la promoción de la reutilización de estas plataformas, el desarrollo de nuevo software y la adaptación parcial de aplicaciones a entornos actuales.

\section{Objetivos}
\label{S:objetivos}

El objetivo de este trabajo es crear un entorno de desarrollo integrado (IDE) en Java que permita el desarrollo de software en C para sistemas de 8 bits. Entre sus características principales se incluyen la compilación cruzada, un sistema de plugins para ampliar funcionalidades y la posibilidad de añadir nuevos compiladores.

Objetivos específicos:

\begin{itemize}
        \item Implementar integración automática con compiladores cruzados.
        \item Diseñar una arquitectura modular basada en plugins.
        \item Permitir la incorporación de nuevas plataformas sin modificar el núcleo del sistema.
        \item Simplificar el flujo de trabajo del desarrollador de software clásico.
    \end{itemize}

El alcance del proyecto se centra principalmente en las plataformas Game Boy y ZX Spectrum, dado que existen librerías ya desarrolladas (GBDK~\cite{gbdk} y Z88DK~\cite{z88dk}) y el software para estas máquinas es de pequeño tamaño, facilitando la realización de pruebas y la validación de la herramienta. Otra gran ventaja de usar estos sistemas es que utilizan arquitecturas basadas total o parcialmente en el procesador Z80 desarrollado por Zilog. Asimismo, se ha diseñado el sistema de plugins para permitir futuras ampliaciones del IDE y facilitar la inclusión de más sistemas desarrollables dentro de la herramienta.

Justificación técnica de la selección de plataformas Gameboy y ZX Spectrum:

\begin{itemize}
        \item Ambas usan CPU basadas en Z80 (total o parcialmente).
        \item Comunidad activa.
        \item Librerías existentes.
        \item Tamaño reducido de ROM → validación rápida.
    \end{itemize}

Los capítulos siguientes presentan la metodología utilizada, el desarrollo de la herramienta, las pruebas realizadas, los resultados obtenidos y las conclusiones derivadas de este proyecto.

\clearpage
\section{Estado del arte}
\label{S:estado_arte}

En este apartado se analizarán los ejemplos más representativos de herramientas para el desarrollo de aplicaciones de 8 bits. Se incluirán tanto IDEs como compiladores. 

\par{La Tabla~\ref{tab:comparativa-ides} recoge un resumen comparativo de las principales características de los entornos de desarrollo más representativos que existen en la actualidad, que serán analizadas en detalle en las subsecciones posteriores. Se prestará especial atención a distinguir entre aquellas funcionalidades integradas de forma nativa y aquellas que dependen de complementos externos o configuraciones adicionales.}

\begin{table}[!h]
    \centering
    \adjustbox{width=\textwidth,center}
    {
        \begin{tabular}{|l|c|c|c|c|c|c|}
            \hline
            \textbf{Característica} & \textbf{GBDK} & \textbf{Z88DK} & \textbf{ZX Basic Studio} & \textbf{VS Code + plugins} & \textbf{Eclipse CDT + toolchains} \\
            \hline
            IDE integrado & No & No & Sí & Sí & Sí\\
            \hline
            Edición de código avanzada & No / mín. & No / mín. & Sí & Sí & Sí\\
            \hline
            Compilación cruzada C & Sí & Sí & No & Sí* & Sí*\\
            \hline
            Soporte múltiples plataformas & No (GB) & Parcial (Z80) & No (ZX Spectrum) & Sí* & Sí* \\ 
            \hline
            Depuración integrada & No & No & Sí & Sí* & Sí* \\
            \hline
            Integración con emuladores & No & No & Sí & Sí* & Sí* \\
            \hline
            Arquitectura de plugins/extensible & No & No & No & Sí* & Sí* \\
            \hline
            Curva de aprendizaje & Alta & Alta & Baja‑media & Media‑alta & Alta \\
            \hline
            Orientado específicamente a 8 bits & Sí & Sí & Sí & Sí* & Sí* \\
            \hline
            Enfoque formativo & Medio & Medio & Medio & Variable & Variable \\
            \hline
            \multicolumn{6}{|l|}{\textit{Nota: El asterisco (*) indica que la característica depende de configuración manual.}} \\
            \hline
        \end{tabular}
    }
    \caption{Comparativa de herramientas y entornos para desarrollo en plataformas de 8 bits}
    \label{tab:comparativa-ides}
\end{table}


%% Seccion 1.3.1 GBDK-2020
\subsection{GBDK-2020}

Como se indica en la documentación oficial, GBDK es un kit de desarrollo cruzado para consolas basadas en SM83, Z80 y 6502. Incluye bibliotecas, utilidades de toolchain y el conjunto de compiladores C SDCC~\cite{gbdk}.

En la actualidad, el proyecto se escindió en una rama diferente llamada GBDK-2020 y el proyecto original, aunque aún es posible encontrarlo en internet, no está mantenido por sus autores Pascal Felber, Lars Malmborg, Michael Hope, David Galloway (djmips), John Fuge y colaboradores.

Incluye tanto las librerías necesarias como el compilador para usarlas además de documentación completa.

Un ejemplo de software desarrollado usando este kit es el proyecto Orange FM, cartucho moderno diseño para la familia de consolas portátiles Game Boy que funciona como un sistema de radio estéreo~\cite{orangefm}. En la Figura~\ref{fig:orangeFM} se muestra la interfaz de este proyecto ejecutándose en la consola Game Boy.

\begin{figure}[!htb]
    \centering
    \includegraphics[width=7cm]{orangeFM.png}
    \caption{Captura de pantalla del proyecto Orange FM.}
    \label{fig:orangeFM}
\end{figure}

Características de GBDK-2020:
\begin{itemize}
    \item Conjunto de herramientas para desarrollo en C y ASM basada en SDCC.
    \item Múltiples bibliotecas y su código fuente.
    \item Programas de ejemplo realizados en ASM y en C.
    \item Soporte para múltiples bancos de memoria ROM.
    \item La mayor parte del código de las herramientas está bajo la licencia GNU GPL. 
    \item GBDK-2020 permite el desarrollo para las siguientes plataformas
    \begin{itemize}
        \item Nintendo Game Boy / Game Boy Color.
        \item Analogue Pocket.
        \item Sega Master System y Game Gear.
        \item Mega Duck / Cougar Boy.
        \item NES.
    \end{itemize}
\end{itemize}

Además incluye diversas utilidades para facilitar los desarrollos:
\begin{itemize}
    \item GBCompress Comprime (y descomprime) datos de archivos binarios con el algoritmo gbcompress.
    \item png2asset Herramienta para convertir archivos PNG a MetaSprites y mapas de mosaicos en formato GBDK.
    \item Romusage Utilidad para estimar el uso de ROM de Game Boy y SMS/GG a partir de archivos .noi y .map, ROM binarias y más.
\end{itemize}

Todo lo anteriormente comentado convierte GBDK-2020 en entorno de desarrollo multiplataforma para sistemas retro de 8 bits, pero sigue dependiendo de un IDE externo o la combinación de editor + terminal para el desarrollo de software.

\clearpage
%% Seccion 1.3.2 Z88DK
\subsection{Z88DK}
Z88DK es un kit de desarrollo cruzado para sistemas basados en Z80. Incluye bibliotecas, utilidades de toolchain y el conjunto de compiladores C SDCC.

Fue creado originalmente por Dominic Morris a finales de la década de 1990 con el objetivo de facilitar el desarrollo de software en lenguaje C para el ordenador Cambridge Z88. Con el tiempo, el proyecto ha evolucionado hasta convertirse en un entorno de desarrollo multiplataforma mantenido por una comunidad de desarrolladores que amplían continuamente el soporte para diferentes sistemas basados en procesadores Z80~\cite{z88dk}.

Como se indica en la documentación oficial: z88dk es el único kit de desarrollo de ensamblador y C que viene listo para usar para crear programas para más de 100 máquinas de la familia z80 (8080, 8085, gbz80, z80, z180, ez80\_z80, KC160, Rabbit 2000, 3000, 4000, 5000)~\cite{z88dk}.

\par{El compilador z88dk soporta un gran número de sistemas basados en procesadores Z80 y arquitecturas relacionadas. La tabla \ref{tab:z88dk_platforms} muestra una selección representativa de algunas de las plataformas más conocidas soportadas por este entorno de desarrollo.}

\par{Algunas de las herramientas incluidas en Z88DK:}
\begin{itemize}
    \item zcces la interfaz de la cadena de herramientas. zcc puede generar un binario de salida a partir de cualquier conjunto de archivos fuente de entrada.
    \item z88dk-scc80 es el compilador C nativo de z88dk.
    \item z88dk-zsdcc es la personalización de z88dk del compilador C optimizado de sdcc.
    \item z88dk-z80asm es un ensamblador/vinculador/bibliotecario con todas las funciones que implementa secciones.
    \item z88dk-zobjcopy permite manipular archivos de objetos y bibliotecas creados por Z80ASM.
    \item z88dk-appmake procesa los binarios generados por el kit de herramientas en un formato adecuado para equipos de destino específicos. Por ejemplo, puede generar archivos hexadecimales de Intel, cintas, ROM, imágenes de disco, etc.
    \item z88dk-lib es un instalador de bibliotecas de terceros. Gestiona la instalación, la eliminación y el listado de bibliotecas disponibles.
    \item z88dk-zx0 y z88dk-zx7 son herramientas de compresión de datos.
\end{itemize}


\par{ ZNAKE \cite{znake_itchio} es un ejemplo de software desarrollado para ZX Spectrum programado en C y compilado con z88dk y compatible con versiones recientes de este compilador, fué publicado en itch.io por marco's retrobits. La Figura~\ref{fig:znake} ilustra una captura de pantalla del juego en ejecución en un emulador de ZX Spectrum.}

\begin{figure}[!htb]
    \centering
    \includegraphics[width=7cm]{znake.png}
    \caption{Captura de pantalla del proyecto ZNAKE.}
    \label{fig:znake}
\end{figure}

\begin{table}[ht]
    \centering
    \small
    \begin{tabular}{p{4cm}p{2cm}p{2.5cm}p{2cm}p{2cm}}
        \toprule
            Machine & Console & Graphics & File I/O & Sound \\
        \midrule
            Altair 8800 & Yes & No & No & No \\
            Amstrad CPC & Yes & 640x200 & Yes & PSG \\
            Apple II & 40x24 & No & CP/M & No \\
            Cambridge Z88 & Yes & 256x64 & Yes & 1 bit \\
            Colecovision & 32x24 & 256x192 & No & PSG \\
            Commodore 128 (Z80) & 40x25 & 640x200 & CP/M & SID + PSG \\
            Enterprise 64/128 & 40x25 & 672x243 & No & Yes \\
            Jupiter Ace & 32x24 & 64x48 & No & 1 bit \\
            MSX & 40x24 & 256x192 & No & PSG \\
            NEC PC-8801 & 40x25 & 600x200 & No & PSG \\
            Nintendo Game Boy & 20x18 & 160x144 & No & Yes \\
            Sega Master System & 32x24 & 256x192 & No & PSG \\
            Sharp X1 & 40x25 & 640x200 & No & PSG \\
            Sinclair ZX80 & 32x24 & 64x48 & No & 1 bit \\
            Sinclair ZX81 & 32x24 & 64x48 & No & 1 bit \\
            ZX Spectrum & 32x24 & 256x192 & Yes & 1 bit + PSG \\
            ZX Spectrum Next & 32x24 & 256x192 & Yes & 1 bit + PSG \\
            SAM Coupé & 32x24 & 512x192 & No & PSG \\
            TRS-80 & 64x16 & 128x48 & No & 1 bit \\
            TI-83+ & 16x8 & 128x64 & No & 1 bit \\
        \bottomrule
    \end{tabular}
    \caption{Selección representativa de plataformas soportadas por el compilador z88dk. Adaptado de la documentación oficial \cite{z88dk_wiki}.}
    \label{tab:z88dk_platforms}
\end{table}


\clearpage

%% 1.3.3 ZX Basic Studio
\subsection{ZX Basic Studio}

\par {ZX Basic Studio \cite{zxbasicstudio_github} es un proyecto en desarrollo, iniciado por un desarrollador que utiliza el seudónimo "Dr. Gusman". Se trata de un IDE multiplataforma orientado al compilador Boriel Basic \cite{boriel_basic_github}, un compilador BASIC de código abierto dirigido a usuarios de ZX Spectrum. Boriel Basic permite escribir programas en BASIC que se compilan a código máquina eficiente, ofreciendo facilidad de uso y rendimiento.}

\par{El IDE ZX Basic Studio está disponible para Windows, Linux y Mac. Integra un emulador y un depurador, así como un editor de gráficos, basado en una versión mejorada de ZX Graphics desarrollada por Juan Segura Duran (“Duefectu”) \cite{duefectu_graphics}. Además, incluye un editor de ficheros .tap que permite gestionar no solo archivos individuales, sino proyectos completos. Su sistema de plugins amplía las funcionalidades del IDE, incorporando editores de gráficos, coloreadores de sintaxis, tooltips de ayuda y otras herramientas de soporte.}

\par{Actualmente, ZX Basic Studio permite compilar y ejecutar programas para ZX Spectrum Next, aunque por el momento la depuración directa en hardware real no está soportada. Dentro del ecosistema ZX Spectrum, ZX Basic Studio se posiciona como una herramienta moderna que combina edición de código, depuración y soporte gráfico, ofreciendo una alternativa más integrada frente a otros IDE tradicionales o emuladores con funcionalidades limitadas. En la Figura~\ref{fig:zxbasicStudio} se puede apreciar la interfaz del IDE con sus diferentes paneles y herramientas integradas.}



\begin{figure}[!htb]
    \centering
    \includegraphics[width=10cm]{zxbasicStudio.png}
    \caption{Captura de pantalla del proyecto ZX Basic Studio.}
    \label{fig:zxbasicStudio}
\end{figure}
    

\clearpage

%% 1.3.4 VS Code + plugins
\subsection{VS Code + plugins}

\par{Visual Studio Code \cite{vscode_history} combina la simplicidad de un editor de código fuente con herramientas avanzadas para desarrolladores, como autocompletado de código y depuración mediante IntelliSense. Su diseño minimalista permite que el flujo de edición, compilación y depuración sea ágil, facilitando la concentración en el desarrollo de software.}

\par{VS Code es un editor rápido, adecuado para un uso intensivo. Ofrece compatibilidad con la mayoría de los lenguajes de programación, incluyendo resaltado de sintaxis, coincidencia de corchetes, sangría automática, selección de columnas y fragmentos de código reutilizables.}

\par{Incluye un depurador interactivo que permite recorrer el código fuente paso a paso, inspeccionar variables, consultar la pila de llamadas y ejecutar comandos en la consola. Además, se integra con múltiples herramientas de compilación y scripting para automatizar tareas rutinarias.}

\par{VS Code es compatible con sistemas de control de versiones como Git, proporcionando un visor de diferencias, historial de cambios y gestión de cambios pendientes. La Figura~\ref{fig:CapturaVSCode} muestra la interfaz principal del editor con un proyecto C abierto.}

\par{Para el desarrollo específico en plataformas de 8 bits, VS Code cuenta con una amplia gama de extensiones y plugins que facilitan el trabajo con estas arquitecturas. La Tabla~\ref{tab:vscode-plugins-8bit} detalla algunas de las extensiones más relevantes disponibles para este tipo de desarrollo.}

\begin{figure}[!htb]
    \centering
    \includegraphics[width=10cm]{CapturaVSCode.png}
    \caption{Captura de pantalla de Visual Studio Code.}
    \label{fig:CapturaVSCode}
\end{figure}


\begin{table}[!h]
    \caption{Plugins de VS Code para desarrollo en plataformas de 8 bits}
    \centering
    \adjustbox{width=\textwidth,center}
    {
        \begin{tabular}{|l|c|c|c|}
            \hline
            \textbf{Plugin / Extensión} & \textbf{Lenguaje / Sistema} & \textbf{Función principal} & \textbf{Notas} \\
            \hline
            Retro Assembler & 6502, Z80, 6809 & Assembler + resaltado de sintaxis & Compatible con NES, C64, ZX Spectrum \\
            \hline
            vscode-asm6502 & 6502 & Resaltado de sintaxis y snippets & Ideal para NES y Commodore 64 \\
            \hline
            z80-code & Z80 & Resaltado de sintaxis y snippets & Compatible con ZX Spectrum, Amstrad CPC \\
            \hline
            GBDK Support & C (GBDK-2020) & Compilación y snippets para Game Boy & Incluye build tasks y ejemplos \\
            \hline
            ZX Basic Studio Integration & Boriel Basic & Soporte de sintaxis y proyectos ZX Basic & Requiere configuración manual de paths \\
            \hline
            C64 Debugger & C64 & Depurador integrado para VICE & Funciona con emulador C64 \\
            \hline
            NES Debugger & 6502 / NES & Depuración de ROMs & Integración con FCEUX \\
            \hline
            Arduino / PlatformIO & C / C++ & Compilación y depuración hardware & Útil para MCU de 8 bits (AVR, 8051) \\
            \hline
            GBDK Snippets & C (Game Boy) & Snippets para librerías GBDK & Acelera desarrollo de gráficos y sonido \\
            \hline
            VIC-20 / C64 Emulation & 6502 & Integración con emuladores & Ejecuta código en entornos simulados \\
            \hline
            Retro Debug & Multi Z80 / 6502 & Depuración general retro & Soporta snapshots y breakpoints \\
            \hline
            ZX Spectrum Next Tasks & Z80 / BASIC & Build y ejecución de programas & Compatible con emuladores y hardware real \\
            \hline
        \end{tabular}
    }
    \label{tab:vscode-plugins-8bit}
\end{table}


\clearpage

%% 1.3.5 Eclipse CDT + toolchains
\subsection{Eclipse CDT + toolchains}

\par{Eclipse es un IDE de código abierto y multiplataforma \cite{eclipse_org}, originalmente diseñado para desarrollo en Java, pero actualmente usado para una amplia variedad de lenguajes y sistemas. Apareció en 2001 como proyecto de IBM y posteriormente pasó a ser administrado por la Eclipse Foundation.}

\par{La Figura~\ref{fig:Eclipse-CDT} muestra la interfaz de Eclipse CDT con un proyecto C configurado, donde se pueden apreciar los diferentes paneles y herramientas disponibles para el desarrollo.}

\begin{figure}[!htb]
    \centering
    \includegraphics[width=10cm]{Eclipse-CDT.png}
    \caption{Captura de pantalla de Eclipse CDT.}
    \label{fig:Eclipse-CDT}
\end{figure}

\par{CDT (C/C++ Development Tooling) \cite{eclipse_cdt} es un conjunto de extensiones que habilita desarrollo en C y C++, incluyendo compilación, depuración, navegación de código, autocompletado y análisis estático. Eclipse CDT es, por tanto, una versión de Eclipse adaptada para proyectos en C. Por sí solo no incluye toolchains, por lo que debe integrarse con compiladores externos como GCC, SDCC, z88dk o GBDK. Esta modularidad permite su uso en plataformas de 8 bits, aunque requiere configuración manual.}

\par{La integración de toolchains externas en Eclipse CDT se realiza mediante la configuración de compiladores y enlazadores dentro del sistema de construcción (Build System). El entorno permite definir rutas a ejecutables, parámetros de compilación, scripts personalizados y perfiles de construcción diferenciados (Debug/Release). En el caso de toolchains como SDCC o z88dk, es necesario definir manualmente los comandos de compilación y enlazado, así como las opciones específicas del microprocesador objetivo.}


\par{Entre sus principales ventajas se encuentran la robustez, estabilidad y amplia compatibilidad con toolchains para 8 bits. Sus principales limitaciones son el alto consumo de recursos y la complejidad de configuración inicial. }


\par{Eclipse CDT surgió en 2005 como plugin del IDE matriz, permitiendo su aplicación en desarrollo de software embebido y retro gracias a la integración con librerías y herramientas de terceros.}


\par{En el contexto específico de desarrollo para sistemas de 8 bits, Eclipse CDT no proporciona asistentes específicos para plataformas retro, por lo que la configuración de entornos como Game Boy, ZX Spectrum o MSX debe realizarse manualmente. Asimismo, la integración con emuladores requiere scripts externos o configuraciones personalizadas.}

\par{La Tabla~\ref{tab:eclipse_cdt} resume las características clave de este entorno de desarrollo para proyectos de 8 bits.}

\begin{table}[!h]
    
    \centering
    \tiny
    \adjustbox{width=0.7\textwidth,center}{
        \begin{tabular}{|l|p{4.5cm}|}
        \hline
        \textbf{Característica} & \textbf{Descripción} \\
        \hline
        IDE completo & Editor, gestión de proyectos y depuración integrada. \\
        \hline
        Integración toolchains & Compatible con GCC, SDCC, z88dk, GBDK. Múltiples plataformas retro. \\
        \hline
        Depuración & Soporte GDB y depuradores específicos de MCU/emuladores. \\
        \hline
        Multiplataforma & Windows, Linux y macOS. \\
        \hline
        Plugins/Extensible & Soporte para nuevos MCU, lenguajes y emuladores. \\
        \hline
        \end{tabular}
    }
    \caption{Características clave de Eclipse CDT para desarrollo en 8 bits}
    \label{tab:eclipse_cdt}
\end{table}


\clearpage

%% Seccion 1.4 Estructura de la memoria
\section{Estructura de la memoria}
\label{S:estructura_memoria}

La memoria del proyecto se organiza en varios capítulos que describen de forma progresiva el contexto, desarrollo y resultados del trabajo realizado:

\begin{enumerate}
    \item \textit{Introducción}: En este apartado se presenta el contexto general del proyecto y la motivación que ha llevado a su desarrollo. Asimismo, se incluye un breve análisis del estado del arte con el fin de situar el trabajo dentro de las soluciones y herramientas existentes.
   
    \item \textit{Análisis}: Se detallan los objetivos del proyecto y las herramientas consideradas para su desarrollo. También se definen los requisitos del sistema y los principales casos de uso que guiarán la implementación.

    \item \textit{Diseño}: Se describe el diseño general del sistema, incluyendo la estructura de las pantallas de la aplicación y el flujo de navegación entre ellas.

    \item \textit{Implementación}: En este capítulo se explican las funcionalidades desarrolladas y se justifican las principales decisiones técnicas adoptadas durante la implementación.

    \item \textit{Pruebas}: Se presenta un conjunto de pruebas destinadas a verificar el correcto funcionamiento del prototipo y a detectar posibles errores en sus funcionalidades principales.

    \item \textit{Planificación y presupuesto}: Se incluye la planificación temporal del proyecto mediante un diagrama de Gantt, así como una estimación de los recursos y costes asociados a su desarrollo.

    \item \textit{Conclusiones y trabajos futuros}: Finalmente, se recogen las conclusiones derivadas del desarrollo del proyecto, junto con una reflexión sobre los conocimientos adquiridos y posibles líneas de mejora o ampliación para trabajos futuros.

\end{enumerate}

\end{document}